\section{Introduction}
\label{sec:ieee-introduction}

Physical AI systems do not act on truth. They act on telemetry, state estimates, and
learned summaries of the world. When that observation channel degrades, the entire
software stack can remain internally coherent while the physical system becomes
externally unsafe. Batteries can violate hidden state-of-charge limits, vehicles can
accelerate on stale headway estimates, and clinical monitoring systems can suppress
necessary intervention. ORIUS is built around the claim that this failure
mode is not a battery problem, an autonomy problem, or a cyber-security problem in
isolation. It is a universal \emph{action-semantics} problem for Physical AI.

This draft therefore asks a narrower but more consequential question than most
controller- or model-centric safety papers: what runtime layer is needed once the
observation surface itself can no longer be trusted to carry actuation legality? The
answer proposed here is ORIUS, Observation--Reality Integrity for Universal Safety, a
typed runtime layer that detects degraded observation, calibrates uncertainty to the
reduced trust surface, constrains admissible actions, repairs or replaces unsafe
candidates, and emits an auditable certificate before release. The contribution is not
one universal policy. It is one universal safety grammar that sits between nominal
intelligence and physical actuation.

The paper is intentionally universal-first. Battery remains the witness row with the
deepest theorem-to-code-to-artifact closure, but it is not treated as the conceptual
center of the argument. Instead, the central object is the observation--action safety
gap (OASG): the event in which an action is legal on the observed state while unsafe on
the true state. That object links conformal and distribution-free uncertainty
\cite{vovk2005algorithmic,romano2019conformalized,gibbs2021adaptive,angelopoulos2023gentle,barber2023beyond},
runtime assurance and supervisory control
\cite{sha2001using,sha1998simplex,desai2019soter,bartocci2018introduction,leucker2009brief},
safety filters and constrained control
\cite{ames2019control,wabersich2021predictive,fisac2019general,ben2009robust},
and anomaly- or degradation-aware cyber-physical monitoring
\cite{rajkumar2010cps,teixeira2015secure,basseville1993detection,pasqualetti2013attack}.
The manuscript's claim is that none of these families alone supplies the runtime bridge
from degraded observation to repaired actuation and certificate semantics across
multiple physical domains.

The technical move is a Detect--Calibrate--Constrain--Shield--Certify kernel. Detect
degraded telemetry and estimate reliability. Calibrate the observation-consistent state
set rather than trusting nominal point estimates. Constrain the action set under that
inflated state uncertainty. Shield the plant by repairing or replacing the candidate
action. Certify the resulting release so fallback, latency, and audit continuity become
part of the scientific object rather than operational afterthoughts. ORIUS makes that
kernel concrete through one adapter contract, one benchmark schema, and one governance
surface. The adapter contract allows the defended Battery, AV, and Healthcare rows to
instantiate the same safety layer without pretending they share the same nominal
controller or safety objective.

The evidence posture of this flagship draft is intentionally visible. ORIUS has one
witness row and two defended bounded rows. The paper therefore argues for a field-level
runtime architecture while keeping the 3-domain promotion boundary explicit in the body
text rather than hiding it in editorial support material.

The paper makes four contributions.
\begin{enumerate}[leftmargin=1.25em,itemsep=0.2em]
\item It defines OASG as a universal safety object for Physical AI and frames degraded
observation as an action-semantics problem rather than only a detection or estimation
problem.
\item It presents ORIUS as a reusable runtime kernel with typed observation,
uncertainty, action-repair, fallback, and certificate objects.
\item It unifies benchmark and governance semantics across the defended 3-domain lane
using one true-state-violation schema and one auditable runtime discipline.
\item It argues, with explicit evidence-tier discipline, that runtime safety layers are
the right universal object for Physical AI safety research even when evidence depth
remains asymmetric across domains.
\end{enumerate}

The rest of the paper follows that logic. It formalizes OASG, positions ORIUS
adversarially against adjacent literatures, details the runtime kernel and the theorem
bridge, walks through the defended Battery, AV, and Healthcare instantiations under one
shared template, and closes with a synthesis that separates universality of contract
from current evidence depth.
